\chapter{Conclusion}
\label{chap:conclusion}

This thesis asked whether a dynamical model of interacting firms can produce the
firm-growth anomalies by itself, without the heavy-tailed internal composition that the
granular models assume and that \citet{moran2024} show the data reject. A Generalized
Lotka--Volterra network whose firms are identical, with no injected
noise and no built-in correlations, recovers the fat-tailed growth-rate distribution, a
size--variance relation that decreases with firm size, as in the data, and every size-conditioned
test by which the data reject the granular picture: the anomalous multiscaling of the
higher volatility moments, the collapse of the rescaled volatility distribution with a
bounded tail, and the persistence of fat tails in the growth of the largest firms. All
come out as stationary properties of an economy that grows while its composition keeps
fluctuating.

The model is one change to the standard Generalized Lotka--Volterra system.
Self-limitation and interactions act on a firm's size relative to the average firm rather
than on its absolute size, which makes the dynamics scale-invariant. Total output then grows
exponentially, without the finite-time blow-up of the absolute model, and in the strongly
disordered regime the composition of the economy never settles but keeps fluctuating
chaotically. The couplings use the standard Lotka--Volterra normalization of the disordered
literature, on a scale-free graph whose mean degree scales with the
number of firms.

The results point to two distinct mechanisms.
Switching the disorder off and on shows the fluctuations and the fat-tailed growth
distribution to be made by the disordered interactions; they appear on every topology,
including a fully connected one. Swapping the degree distribution of the graph, with
everything else held fixed, shows the size-conditioned structure (the low
size--variance exponent and the multiscaling) to be made by the degree heterogeneity.
Under the standard normalization a firm's noise variance grows with its connectivity, so
the graph's degree distribution sets the distribution of noise levels across firms. On
homogeneous graphs the multiscaling vanishes exactly; as the degree distribution
broadens the ratio curve bends toward the data, and only the scale-free graph reaches
it. The same heterogeneity reaches the individual firm through survival: most highly
connected firms are driven to the immigration floor, and the survivors'
sizes grow with degree as $\bar S \sim k^{0.56}$, so the degree hierarchy shapes
the size hierarchy as well as the noise.
A network account of the \emph{level} of firm volatility, resting on customer
concentration, already exists \citep{herskovic2020}; the present model adds that the
\emph{size-conditioned} statistics follow from the heterogeneity of the interaction
graph.

Three limitations are important. First, neither the dynamics nor the construction is new,
and Chapter~\ref{chap:background} places both in their lineage: the fluctuating
regime is the chaotic phase of the disordered Lotka--Volterra model \citep{diederich1989,bunin2017,galla2018,roy2020},
and the scale invariance and mean-coupling are those of the econophysics models of wealth and
firm sizes \citep{bouchaudmezard2000,malcai2002}. What is
specific here is the combination, self-limitation and disordered interactions made relative
together, and the finding that the graph's heterogeneity generates the discriminating
firm-growth statistics. Second,
the quantitative match is partial. The
multiscaling ratios and the conditional tests (the shape of the size-conditioned
statistics) are reproduced and are converged in system size. The absolute scale is
not on the same graph: on the power-law-$2.5$ graph of the headline figures the
size--variance exponent stands at $\beta \approx 0.45$ at $N = 16000$ and
$0.49$ at $N = 32000$,
against the empirical $0.15$--$0.20$, still drifting slowly upward, while on the
power-law-$2.2$ graph the exponent is converged at the upper edge of the empirical band
($\beta \approx 0.21$--$0.23$) with the multiscaling ratios then running
$5$--$7\%$ below the data; no single tail exponent yet places both on the data at
once. The model robustly predicts the shape of the statistics, while their scale remains a finite-size
quantity. Third,
the growth is near-critical: the operating point sits in the growing band but not far
from its boundaries; stronger disorder loses realizations to collapse, and a
lighter-tailed graph pushes them into freezing. Perturbing the operating point
($\sigma$ from $1.5$ to $2.0$, $\mu$ from $-1.4$ to $-2.1$) moves the size--variance
exponent substantially ($\beta = 0.28$--$0.79$) while the ratio curve stays bent below
the fixed-shape line wherever the fluctuating state survives, so the multiscaling is a
property of the phase and the particular exponent values of the operating point.

\enlargethispage{\baselineskip}
The exponent's slow upward drift calls for determining its limiting value and for developing an analytical
theory. The mean-field treatment of Appendix~\ref{app:dmft} stops at the relaxed phase
on the homogeneous graph. The
size distribution, measured in Appendix~\ref{app:size-distribution}, is effectively
truncated where the data are Zipf-like, so matching it will take a modified model, not a
larger simulation.
The immigration floor's plateau on the smallest firms is worth a second look,
now that \citet{fontanelli2025} report an empirical size--volatility relation
that steepens for small firms and flattens for large ones.
The correlation structure is the sharper question. \citet{moran2024} conjecture that
the granular models fail for lack of a between-firm growth correlation that grows
with size; Appendix~\ref{app:correlations} measures that correlation directly and finds it modest
and roughly independent of both size and degree. The model therefore reproduces the
size-conditioned statistics through the degree-dependence of the noise amplitude, not
through the size-dependent correlations proposed by Moran et al. Whether real firms
carry the size-growing or the flat correlation is a question for firm-level data on a
known network. Incorporating couplings from real input--output and supply linkages would connect the model
more directly to the macroeconomic dynamics behind these stylized facts.
